% GENERATED FILE. Edit canonical lessons, then run npm run generate:books.
% canonical-source-hash: fnv1a64:ccd51561bca3771d
% canonical-lessons: BN-C257-chhana, BN-C257-ponir, BN-C257-aiskrim, BN-C257-chire, BN-C257-phuchka

\chapter{Chhana, Paneer, Ice cream, Chira, Phuchka}
\label{ch:bn-chhana-paneer-ice-cream-chira-phuchka}
\hlchaptermodality{257}

\begin{chapteropening}
\textbf{By the end of this chapter:} \emph{I can say chhana, paneer, ice cream, chira and phuchka.}

Complete the last lesson of chapter 257: I can say chhana, paneer, ice cream, chira and phuchka.
\end{chapteropening}

\section[\texorpdfstring{\bn{ছানা} (chhānā)}{ছানা (chhānā)}]{\bn{ছানা} (chhānā) — chhana (fresh cheese)}

\label{lesson:BN-C257-chhana}



\begin{quote}
Before the new one: say the Bengali for art; industry, then the Bengali for literature.
\end{quote}

\subsection*{You'll want to know: \bn{ছানা}}
\textbf{\bn{ছানা}} — \emph{chhānā} — \textquotedblleft{}chhana (fresh cheese)\textquotedblright{}.

\begin{grammarlens}[title={What you've built}]
One more word for this chapter.
\end{grammarlens}

\subsection*{Your turn}
\begin{itemize}
\raggedright
  \item \emph{Say it:} \emph{chhānā}
  \item \emph{Say it:} \emph{chhānā}, once more
  \item \emph{Say it:} say \emph{chhānā}
  \item \emph{Recall:} say the Bengali for art; industry, then the Bengali for literature, then say \emph{chhānā} again
\end{itemize}

\begin{tcolorbox}[breakable,colback=teal!4,colframe=teal!35!black,title={Before you move on}]
What does \bn{ছানা} mean? (\textquotedblleft{}Chhana (fresh cheese)\textquotedblright{}.) Say it once more.
\end{tcolorbox}
\section[\texorpdfstring{\bn{পনির} (pônir)}{পনির (pônir)}]{\bn{পনির} (pônir) — paneer}

\label{lesson:BN-C257-ponir}



\begin{quote}
Before the new one: say the Bengali for literature, then the Bengali for chhana.
\end{quote}

\subsection*{You'll want to know: \bn{পনির}}
\textbf{\bn{পনির}} — \emph{pônir} — \textquotedblleft{}paneer\textquotedblright{}.

\begin{grammarlens}[title={What you've built}]
One more word for this chapter.
\end{grammarlens}

\subsection*{Your turn}
\begin{itemize}
\raggedright
  \item \emph{Say it:} \emph{pônir}
  \item \emph{Say it:} \emph{pônir}, once more
  \item \emph{Say it:} say \emph{pônir}
  \item \emph{Recall:} say the Bengali for literature, then the Bengali for chhana, then say \emph{pônir} again
\end{itemize}

\begin{tcolorbox}[breakable,colback=teal!4,colframe=teal!35!black,title={Before you move on}]
What does \bn{পনির} mean? (\textquotedblleft{}Paneer\textquotedblright{}.) Say it once more.
\end{tcolorbox}
\section[\texorpdfstring{\bn{আইসক্রিম} (āiskrim)}{আইসক্রিম (āiskrim)}]{\bn{আইসক্রিম} (āiskrim) — ice cream}

\label{lesson:BN-C257-aiskrim}



\begin{quote}
Before the new one: say the Bengali for chhana, then the Bengali for paneer.
\end{quote}

\subsection*{You'll want to know: \bn{আইসক্রিম}}
\textbf{\bn{আইসক্রিম}} — \emph{āiskrim} — \textquotedblleft{}ice cream\textquotedblright{}.

\begin{grammarlens}[title={What you've built}]
One more word for this chapter.
\end{grammarlens}

\subsection*{Your turn}
\begin{itemize}
\raggedright
  \item \emph{Say it:} \emph{āiskrim}
  \item \emph{Say it:} \emph{āiskrim}, once more
  \item \emph{Say it:} say \emph{āiskrim}
  \item \emph{Recall:} say the Bengali for chhana, then the Bengali for paneer, then say \emph{āiskrim} again
\end{itemize}

\begin{tcolorbox}[breakable,colback=teal!4,colframe=teal!35!black,title={Before you move on}]
What does \bn{আইসক্রিম} mean? (\textquotedblleft{}Ice cream\textquotedblright{}.) Say it once more.
\end{tcolorbox}
\section[\texorpdfstring{\bn{চিঁড়ে} (chĩṛe)}{চিঁড়ে (chĩṛe)}]{\bn{চিঁড়ে} (chĩṛe) — chira (flattened rice)}

\label{lesson:BN-C257-chire}



\begin{quote}
Before the new one: say the Bengali for paneer, then the Bengali for ice cream.
\end{quote}

\subsection*{You'll want to know: \bn{চিঁড়ে}}
\textbf{\bn{চিঁড়ে}} — \emph{chĩṛe} — \textquotedblleft{}chira (flattened rice)\textquotedblright{}.

\begin{grammarlens}[title={What you've built}]
One more word for this chapter.
\end{grammarlens}

\subsection*{Your turn}
\begin{itemize}
\raggedright
  \item \emph{Say it:} \emph{chĩṛe}
  \item \emph{Say it:} \emph{chĩṛe}, once more
  \item \emph{Say it:} say \emph{chĩṛe}
  \item \emph{Recall:} say the Bengali for paneer, then the Bengali for ice cream, then say \emph{chĩṛe} again
\end{itemize}

\begin{tcolorbox}[breakable,colback=teal!4,colframe=teal!35!black,title={Before you move on}]
What does \bn{চিঁড়ে} mean? (\textquotedblleft{}Chira (flattened rice)\textquotedblright{}.) Say it once more.
\end{tcolorbox}
\section[\texorpdfstring{\bn{ফুচকা} (phuchkā)}{ফুচকা (phuchkā)}]{\bn{ফুচকা} (phuchkā) — phuchka}

\label{lesson:BN-C257-phuchka}



\begin{quote}
Before the new one: say the Bengali for ice cream, then the Bengali for chira.
\end{quote}

\subsection*{You'll want to know: \bn{ফুচকা}}
\textbf{\bn{ফুচকা}} — \emph{phuchkā} — \textquotedblleft{}phuchka\textquotedblright{}.

\begin{grammarlens}[title={What you've built}]
One more word for this chapter.
\end{grammarlens}

\subsection*{Your turn}
\begin{itemize}
\raggedright
  \item \emph{Say it:} \emph{phuchkā}
  \item \emph{Say it:} \emph{phuchkā}, once more
  \item \emph{Say it:} say \emph{phuchkā}
  \item \emph{Recall:} say the Bengali for ice cream, then the Bengali for chira, then say \emph{phuchkā} again
\end{itemize}

\begin{tcolorbox}[breakable,colback=teal!4,colframe=teal!35!black,title={Before you move on}]
What does \bn{ফুচকা} mean? (\textquotedblleft{}Phuchka\textquotedblright{}.) Say it once more.
\end{tcolorbox}
